\section{Methods}
\label{sec:methods}

\placeholder{Venue decides the form: STAR Methods (Cell Systems) or a Methods section
after the Discussion (PLOS CB). Content is the same.}

\subsection{The hybrid driver}
\label{sec:driver}

The driver keeps two blocks with separate state: metabolite pools in mM as floating
point values, and molecule counts as integers. One handshake step does the
following, in order:
\begin{enumerate}
  \item update growth and cell geometry from the membrane protein counts;
  \item write the protein counts into the metabolic rate laws (catalytic edges);
  \item step the ODE block by 1\,s (Rodas5P, absolute tolerance $10^{-10}$, relative
  tolerance $10^{-8}$);
  \item debit the pools by the costs of the stochastic events of the previous step
  (deferred-cost edges), clipped at zero as in the published model;
  \item every 60\,s, rebuild the stochastic rate constants from the pools;
  \item step the jump block by 1\,s with Gillespie's direct method
  \citep{gillespie1977exact}.
\end{enumerate}
On the jump clock the rebuilt constants therefore take effect at $60m - 1$\,s. Pools
are written back as molecule counts under a declared rounding policy. A recorded path
can be replayed: the replay applies the recorded firings instead of drawing new ones,
and it returns the log density of the path and its sufficient statistics (event
counts and exposures).

\subsection{Observation model}
\label{sec:obs}

Each synthetic dataset has $C$ cells (\pn{mzero.cells} on M0), observed every 60\,s.
\begin{itemize}
  \item \textbf{Transcripts:} the exact count of each gene's transcripts in each cell
  at each save.
  \item \textbf{Bulk metabolites:} one measurement per pool per save of the mean over
  the $C$ cells, $\bar{x}$, with lognormal noise: $\log z \sim
  \mathcal{N}(\log \bar{x}, \sigb^2)$. The assay noise $\sigb$ is a free parameter.
\end{itemize}
We use bulk metabolites because a per-cell metabolite panel cannot be measured.
\placeholder{State the panel (M0\_PANEL) and the prior on $\sigb$.}

\subsection{The exact sampler}
\label{sec:sampler}

The sampler is a three-block Gibbs scan with the stochastic path as a latent variable.
\begin{enumerate}
  \item \textbf{Paths.} Conditional sequential Monte Carlo over 60\,s windows
  \citep{andrieu2010particle}, with particle Gibbs at $N = 20$ particles. Each
  particle is a snapshot of the driver. Transcripts are proposed as exact
  endpoint-conditioned birth--death bridges \citep{hobolth2009simulation}, so every
  proposal agrees with the observed transcript counts. The other reactions are then
  simulated given the transcripts. The weight is the bulk metabolite likelihood
  times each gene's transition probability. Cells are updated in turn, because the
  bulk likelihood couples them.
  \item \textbf{Metabolic constants and $\sigb$.} A univariate slice step on each log
  parameter, given the paths. Reverse constants follow from the forward ones and the
  equilibrium constants.
  \item \textbf{Gene-expression constants.} Given the path, each promoter strength
  and the decay rate have an exact one-dimensional conditional of the form
  $n u - e^{u} A + \log p(u)$, where $n$ is the event count and $A$ the exposure. We
  sample it by a slice step.
\end{enumerate}
The sampler is written directly, not in a probabilistic programming language. The
path update needs ancestor sampling and exact snapshots of the driver, and a general
framework could not provide them without sharing solver state between particles.
Runs can be checkpointed and resumed bit for bit, and the threaded sweep is bit for
bit equal to the serial one.

\subsection{The transcript factor}
\label{sec:txfactor}

When the paths are integrated out, each gene's transcripts form a birth--death chain
with a piecewise-constant birth rate, observed exactly every 60\,s. The likelihood of
one cell's record is a product of transition probabilities. We compute them by
uniformisation with tridiagonal products, truncated where the remaining mass is below
$10^{-12}$. The cap on the count is fixed per dataset, and the number of terms comes
from the primal value, so the factor is differentiable in the rates.

\subsection{The mean-field route}
\label{sec:meanfield}

The swap replaces the three jump modules with one mean-field module. Its equations
are generated from the jump modules' own propensities and update rules. The composed
model is one ODE problem. Four edges change form, and each change is a declared
reduction:
\begin{itemize}
  \item catalytic edges read the mean protein counts;
  \item drains are continuous and not clipped, and the solve stops if a pool becomes
  negative;
  \item the rebuild is a discrete callback at $60m - 1$\,s and at $t = 0$;
  \item growth is continuous.
\end{itemize}
The bulk factor is $\mathcal{N}\!\left(\log z;\ \log \max(\bar{x}_{\mathrm{MF}}(\theta), 1)
+ \delta,\ \sigb^2 I + S\right)$, with $\delta$ and $S$ from an ensemble of the exact
hybrid at a reference parameter value, and the floor of one molecule.
\placeholder{B5: state how S is estimated (delta method over 800 cells) and the
sampler settings once built.}

\subsection{Simulation-based calibration}
\label{sec:sbc}

For each of $R = \pn{vseven.R}$ replicates we drew a parameter set from the prior,
simulated a dataset, and ran one chain of 656 sweeps from a prior draw. We dropped
100 sweeps and thinned the rest to $L = 100$ draws. The rank of the truth among the
draws is uniform if the sampler is calibrated \citep{talts2018validating}. We tested
each parameter's rank-ECDF against a simultaneous band at $1 - 0.05/6$
\citep{sailynoja2022graphical}, so the test as a whole is at the 95\% level. The pass
rule was fixed before any rank was computed.

\subsection{Compute}
\label{sec:compute}

All runs used the OzSTAR cluster. \placeholder{Hardware, total CPU-hours, and a table
of every job behind each number, from the result files.}
